\section{Grid overlaps, rotations, and full-reconstruction limits}
\label{app:grids}

These calculations arose during the discussion. They delimit what weighted
overlap alone can guarantee and explain why full serialization is not the main
proposal.

\subsection{Independent real-valued grids}

Flatten each grid to \(X_\ell\in\R^{81}\). For
\[
 Y=\sum_{\ell=0}^{K-1}w_\ell X_\ell,
\]
the measurement matrix is
\[
 M=[\,w_0I_{81}\ \cdots\ w_{K-1}I_{81}\,].
\]
If at least one weight is nonzero, \(\rank M=81\), hence
\(\dim\ker M=81(K-1)\). For two nonzero weights, explicit collisions are
\[
 X_i'=X_i+w_jT,\qquad X_j'=X_j-w_iT.
\]
The changes cancel in \(Y\). A sufficiently small nonzero \(T\) also gives
collisions in the interior of bounded value ranges. Different weights cannot
remove this linear ambiguity. A linear representation of \(K\) independent
continuous grids needs at least \(K\) scalar-grid output channels.

\subsection{A weighted shadow with retained contrasts}

For positive weights, retain
\[
 S=\sum_{\ell=0}^{K-1}w_\ell X_\ell,\qquad
 D_\ell=X_\ell-X_0,\quad 1\le\ell<K.
\]
Let \(W=\sum_\ell w_\ell\). Substituting \(X_\ell=X_0+D_\ell\) gives
\[
 S=WX_0+\sum_{\ell=1}^{K-1}w_\ell D_\ell,
\]
so the inverse is
\[
 X_0=\frac{S-\sum_{\ell=1}^{K-1}w_\ell D_\ell}{W},
 \qquad X_\ell=X_0+D_\ell.
\]
There are \(K\) channels, and the signed contrasts must not be clipped.

A well-conditioned alternative is \(y=Qx\) at each cell, where \(Q\) is
orthogonal and its first row is \(w^\top/\norm w_2\). Then
\[
 x=Q^\top y,\qquad \norm{Qx}_2=\norm x_2,\qquad \kappa_2(Q)=1.
\]
For \(a=w/\norm w_2\ne e_0\), an explicit completion is
\[
 v=e_0-a,\qquad Q=I-\frac{2vv^\top}{v^\top v}.
\]
Direct multiplication gives \(Q^\top=Q\), \(Q^2=I\), and \(Qe_0=a\),
so its first row is \(a^\top\). If \(a=e_0\), use \(Q=I\).
These are exact-real constructions, not bitwise floating-point guarantees.

\subsection{Known quarter-turn copies of one source}

Let \(R\) be the exact \(90^\circ\) rotation permutation of the grid, so
\(R^4=I\). For ordered layers \(X,RX,R^2X,R^3X\) and increasing weights
\(1,2,4,8\), set
\begin{equation}
 Y=\frac{X+2RX+4R^2X+8R^3X}{15}.
 \label{eq:rotation}
\end{equation}
Multiplication yields
\[
 (2R-I)(I+2R+4R^2+8R^3)=16R^4-I=15I,
\]
and therefore
\begin{equation}
 \boxed{X=2RY-Y.}
\end{equation}
The descending convention
\((8X+4RX+2R^2X+R^3X)/15\) instead has inverse \(X=2Y-RY\).
Coefficient order and inverse must remain paired.

For any \(r>1\),
\[
 T_r=\frac{I+rR+r^2R^2+r^3R^3}{1+r+r^2+r^3},
 \qquad
 T_r^{-1}=\frac{rR-I}{r-1}.
\]
Indeed,
\((rR-I)(I+rR+r^2R^2+r^3R^3)=(r^4-1)I\),
and \(r^4-1=(r-1)(1+r+r^2+r^3)\).

A \(9\times9\) grid has one fixed center and twenty four-cell orbits.
The rotation eigenvalues \(1,-1,i,-i\) therefore have multiplicities
\(21,20,20,20\). Since \(R\) is an orthogonal permutation, a polynomial
in \(R\) is normal, and its singular values are the magnitudes of its
eigenvalues. At \(r=2\) they are
\[
 1\ (21\text{ times}),\qquad
 1/\sqrt5\ (40\text{ times}),\qquad
 1/3\ (20\text{ times}).
\]
Thus the condition number is \(3\), and
\(\norm{(2R-I)\varepsilon}_2\le3\norm{\varepsilon}_2\).
The same factor bounds maximum cell error by the triangle inequality.
More generally,
\(\kappa_2(T_r)=(r+1)/(r-1)\), because the largest and smallest
inverse magnitudes occur at the rotation eigenvalues \(-1\) and \(1\).

Equal averaging \((I+R+R^2+R^3)/4\) retains only one average per orbit,
so it has rank \(21\) and nullity \(60\). For arbitrary four weights, the
map is invertible on this grid exactly when
\[
 p(t)=w_0+w_1t+w_2t^2+w_3t^3
\]
is nonzero at \(1,-1,i,-i\). Strictly increasing positive weights satisfy
this: \(p(1)>0\), \(p(-1)=(w_0-w_1)+(w_2-w_3)<0\), and
\[
 p(i)=(w_0-w_2)+i(w_1-w_3)\ne0.
\]
They can still be poorly conditioned as they approach equal weighting.

For original eight-bit cells, store the unnormalized integer
\[
 N=X+2RX+4R^2X+8R^3X.
\]
Its cells lie in \([0,3825]\), so 12 bits suffice per output cell.
The inverse \(X=(2RN-N)/15\) is exact on valid codes with sufficiently
wide integer intermediates; each numerator is divisible by 15.

There is only one source grid's worth of independent unknowns here.
The result cannot encode four independent real-valued grids into one.
Quarter-turn permutations, known orientation convention, and known transform
order are assumptions. Arbitrary-angle interpolation or cropping needs
separate analysis. The map is rotation-equivariant, \(T_rR=RT_r\), rather
than rotation-invariant. Discarding orientation while asking to recover it
would require additional pose information.

\subsection{Exact positional packing for a finite alphabet}

For original \(b\)-bit unsigned values, let \(q=2^b\) and
\(x_\ell\in\{0,\ldots,q-1\}\). Define
\begin{equation}
 C=\sum_{\ell=0}^{K-1}q^\ell x_\ell,\qquad
 x_\ell=\left\lfloor C/q^\ell\right\rfloor\bmod q.
 \label{eq:radix}
\end{equation}
The contribution from the highest differing position \(j\) differs by
at least \(q^j\). The greatest possible opposing lower contribution is
\[
 (q-1)\sum_{i<j}q^i=q^j-1.
\]
It cannot cancel the higher difference. Division and reduction modulo \(q\)
then give the stated inverse.

For arbitrary scalar weights, injectivity on this alphabet is equivalent to
\[
 \sum_\ell w_\ell\delta_\ell\ne0
 \quad\text{for every nonzero integer }\delta
 \text{ with }|\delta_\ell|\le q-1.
\]
A collision gives such a difference vector. Conversely, every such difference
is realized by \(x_\ell=\max(\delta_\ell,0)\) and
\(y_\ell=\max(-\delta_\ell,0)\).
A sufficient dominance condition is
\[
 w_j>(q-1)\sum_{i<j}w_i,\qquad w_0>0.
\]
The highest-position argument proves it. It is not necessary: binary
weights \(2,3,4\) have eight distinct subset sums even though the last
weight does not exceed the sum of the earlier two.

Represent a variable-depth ordered stack by \((C,K)\). Placing stack \(B\)
above \(A\) gives
\[
 (C_A,K_A)\star(C_B,K_B)
 =\bigl(C_A+q^{K_A}C_B,\ K_A+K_B\bigr).
\]
Both groupings of three stacks have code
\[
 C_A+q^{K_A}C_B+q^{K_A+K_B}C_C,
\]
which proves associativity. Given the lower split depth, reduction modulo
\(q^{K_A}\) and integer division recover the two blocks.
Appending a new upper layer uses \(C_{\rm new}=C+q^Kx_{\rm new}\);
the update \(qC+x_{\rm new}\) instead makes older layers heavier.

Depth must be retained because a zero upper layer changes no value code.
Grouping boundaries and source identities require metadata if they count as
information. Values outside the digit range create carries:
\((q,0)\) and \((0,1)\) encode identically. Known unequal digit ranges
admit mixed-radix weights \(w_\ell=\prod_{j<\ell}q_j\), with their bases
available to the decoder.

\subsection{Capacity and brightness precision}

There are \(q^K=2^{Kb}\) independent cell stacks. A fixed-length lossless code
therefore needs at least \(Kb\) bits per cell, or \(81Kb\) bits for the grid,
excluding metadata. A single spatial grid can carry the stack only by
increasing precision or retaining additional state. Positional packing meets
this payload count exactly.

\Needspace{12\baselineskip}
For \(K\ge1\), normalize the integer code to \(z=C/(q^K-1)\), and write
\(M=q^K-1\). Adjacent valid brightness levels are separated by
\(\Delta=1/M\). With exact scaling and rounding, decoding by
\[
 \widehat C=\operatorname{round}(M\widehat z)
\]
is guaranteed when
\[
 |\widehat z-z|<\frac1{2M}.
\]
If the implemented scaled value is \(M\widehat z+\delta\), a sufficient
condition including that arithmetic error is
\[
 M|\widehat z-z|+|\delta|<\frac12.
\]
For four eight-bit layers and exact decoding arithmetic, the normalized
error threshold is about \(1.16415\times10^{-10}\).
For any \(N\) distinct scalar codes in \([0,1]\), the \(N-1\) gaps sum to
at most one, so the minimum gap is at most \(1/(N-1)\).
Alternative weights cannot remove that fixed-range robustness limit.

This is a full-reconstruction reference and a capacity check.
Quantizing original floating-point values into new eight-bit shades is not
lossless. Preserving original floating-point bit patterns requires their full
stored bits, including sign and exponent. Such packing does not automatically
give a useful numerical representation for learning.

\section{ReLU, reversible computation, and numerical exactness}
\label{app:arithmetic}

For real scalars, \((\rho(x),\rho(-x))\) recovers \(x\).
A related bound applies to one bias-free layer
\(h(x)=\rho(Ax)\) on all of \(\R^n\): global injectivity requires at least
\(2n\) output coordinates \cite{injective}. Ignore identically zero rows
and choose a generic point away from the hyperplanes of the nonzero rows.
Active rows must span \(n\) dimensions for local injectivity.
At \(-x\), the complementary nonzero rows must also span \(n\).
There must therefore be at least \(2n\) rows. Pairing the identity and
negative identity attains the bound. It does not impose the same requirement
on every architecture, restricted domain, or target-specific encoding.

\Needspace{12\baselineskip}
ReLU can occur inside a larger invertible operation. For example,
\[
 u=A+F(B),\qquad v=B+G(u)
\]
has inverse
\[
 B=v-G(u),\qquad A=u-F(B).
\]
The functions \(F,G\) themselves need not be invertible. Additive coupling is
a standard construction in invertible networks \cite{nice}. Its full carried
state retains the dimensions needed for the inverse; it does not compress
arbitrary new independent inputs into insufficient capacity.

Algebraic reversibility is distinct from bitwise reversibility. IEEE float32
has 24 significand bits and float64 has 53 \cite{ieee}. With ordinary
round-to-nearest,
\[
 \operatorname{float32}(2^{24}+1)
 =\operatorname{float32}(2^{24}),
\]
and the analogous binary64 collision occurs at \(2^{53}+1\).
A floating-point additive coupling can therefore map different inputs to the
same stored output even though its real-arithmetic formula has an inverse.

For unnormalized integer-valued radix codes, all \(Kb\)-bit payloads are
representable in float32 if \(Kb\le24\), and in float64 if \(Kb\le53\),
with sufficient exponent range and exact intermediate operations.
Correctly implemented unsigned \(Kb\)-bit integer storage retains the entire
payload. Original float32 payloads require 32 stored bits each, not merely
their 24-bit significands.

For fixed-width integer states and deterministic integer-valued maps
\(F,G\) on the appropriate finite state spaces, modular coupling
\[
 u=(A+F(B))\bmod M,\qquad v=(B+G(u))\bmod M
\]
has an exact inverse by modular subtraction. Integer bijective flows provide
a related discrete approach \cite{idf}. Wraparound changes the monotone
darkness interpretation, and training a discrete transformation needs its own
analysis.

Report three different claims separately: algebraic target sufficiency,
stability at a stated error level, and bitwise recovery when actually required.
The selective architecture is not granted the last claim by the first.
